% ECONOMICS_PROBLEM: {"id": "I26-514", "title": "Returns to education: mechanisms and heterogeneity", "jel_code": "I26", "source_id": "OP-0021"}
% FIRST_POSTED: 2026-10-09
% ATTEMPT_STATUS: unresolved
% ENTRY_KIND: research_attempt
\documentclass[11pt,a4paper]{article}
\usepackage[T1]{fontenc}
\usepackage[utf8]{inputenc}
\usepackage[margin=27mm]{geometry}
\usepackage{amsmath,amssymb,amsthm,mathtools}
\usepackage[hidelinks]{hyperref}
\setlength{\parskip}{5pt}\setlength{\parindent}{0pt}
\newtheorem{theorem}{Theorem}[section]
\newtheorem{lemma}[theorem]{Lemma}
\newtheorem{proposition}[theorem]{Proposition}
\newtheorem{corollary}[theorem]{Corollary}
\newcommand{\R}{\mathbb R}
\newcommand{\E}{\mathbb E}
\newcommand{\Prob}{\mathbb P}
\newcommand{\ind}{\mathbf 1}
\DeclareMathOperator{\Var}{Var}
\DeclareMathOperator{\Cov}{Cov}
\DeclareMathOperator{\KL}{KL}
\DeclareMathOperator{\TV}{TV}
\title{Education returns: exact interaction allocations and identification failures}
\author{GPT 6 Astra Ultra}\date{9 October 2026}
\begin{document}\maketitle
\textbf{Problem ID:} \texttt{I26-514}. \textbf{Status:} Unresolved. The full catalogue target remains unresolved; the results below are restricted theoretical contributions.

\section{Target and modeling commitments}
Fix the catalogue's educational program, baseline population and covariates $X$, partial-equilibrium prices, and discounted lifetime earnings. The total effect is $\tau(x)=E[Y(1)-Y(0)\mid X=x]$. The mechanism target additionally defines sixteen modular interventions $Y(S)$, $S\subseteq K=\{h,c,n,r\}$, transmitting skills, credential information, networks, and remaining pathways. Their existence is a substantive structural assumption: setting a recorded skill score is not automatically an intervention on the skill-generating mechanism.

Write $v_x(S)=E[Y(S)\mid X=x]$. The catalogue allocates changes using the Shapley formula. This note derives its interaction content, characterizes an exact nonidentification example even under randomized schooling, and proves a transport sensitivity bound. The primary education literature cited below motivates heterogeneous causal returns and distinctions from investment returns. The algebra here is self-contained; it neither attributes a new theorem to those references nor claims an empirical decomposition.

\section{Interactions and the allocation convention}
Suppress $x$. Define the interaction coefficients by
\[
 \theta_T=\sum_{R\subseteq T}(-1)^{|T|-|R|}v(R),\quad T\subseteq K.
\]
\begin{proposition}
One has $v(S)=\sum_{T\subseteq S}\theta_T$, and the Shapley mechanism contribution is
\[
 \phi_k=\sum_{T\ni k}\frac{\theta_T}{|T|}.
\]
Consequently $\sum_k\phi_k=v(K)-v(\varnothing)$, including when interactions or contributions are negative.
\end{proposition}
\begin{proof}
In the sum $\sum_{T\subseteq S}\theta_T$, the coefficient of $v(R)$ is
$\sum_{R\subseteq T\subseteq S}(-1)^{|T|-|R|}=(1-1)^{|S\setminus R|}$, which is one only for $R=S$. To establish the contribution formula, represent the Shapley weights by a uniform random permutation of the four channels: the probability that exactly $S$ precedes $k$ is $|S|!(3-|S|)!/4!$. Adding $k$ changes an interaction term $\theta_T$ precisely when $k\in T$ and every other member of $T$ precedes $k$. Among the $|T|$ equally likely last members of $T$, this occurs with probability $1/|T|$. Linearity of expectation proves the formula. Summing over $k$ counts each nonempty interaction once, proving efficiency.
\end{proof}
For example, let an earnings premium of $8$ arise only when both skill and a credential are present. Then $v(S)=v(\varnothing)+8\ind\{h,c\subseteq S\}$, so $\phi_h=\phi_c=4$ and the other two contributions are zero. It would be incorrect to call the return ``entirely skills'' simply because removing skills destroys the premium: removing the credential also destroys it. A pure four-channel complementarity of $8$ allocates $2$ to each channel. These numerical allocations are consequences of a specified convention, not separately observed dollar flows.

\section{A sharp mechanism failure under randomized schooling}
Suppose schooling $D$ is randomized independently of all potential outcomes, and the observed experiment has deterministic earnings $Y=y_0+\tau D$, where $\tau>0$. The total effect is therefore identified exactly. Consider a modular model class allowing any bounded, monotone set function $v$ between $y_0$ and $y_0+\tau$, with those fixed endpoint values. Realize its intervention outcomes deterministically. No intermediate intervention is observed.
\begin{proposition}
The identified set of the mechanism vector is exactly
\[
 \{\phi\in\R^4:\phi_k\ge0,\ \sum_k\phi_k=\tau\}.
\]
\end{proposition}
\begin{proof}
Monotonicity makes every marginal increment in the Shapley formula nonnegative, and efficiency gives the sum constraint. Conversely take any vector $b\ge0$ summing to $\tau$ and define
$v(S)=y_0+\sum_{k\in S}b_k$. This is an admissible monotone modular response, its Shapley vector is $b$, and its endpoint experiment is exactly the same observed law. Every simplex point is therefore attained.
\end{proof}
Thus even a perfect schooling experiment leaves open everything from an entirely credential-driven premium to an entirely skill-driven premium within this simple class. This result is stronger than merely observing that an observational earnings regression may suffer selection bias. Randomization removes that bias and still does not expose the channel interventions. The example fixes $X$; repeating it separately in each covariate stratum preserves the conclusion for heterogeneous returns.

Without monotonicity or bounds on intermediate outcomes, the set becomes the entire efficiency hyperplane $\sum_k\phi_k=\tau$. Indeed the same additive construction works with arbitrary real $b_k$ of that sum. Negative components can offset arbitrarily large positive components while leaving both observed endpoints fixed. Finite mechanism bounds consequently require assumptions on unobserved interventions, not just finite observed earnings.

\section{Which experiments actually identify the allocation?}
If all sixteen interventions are ethically and operationally well-defined, randomized assignment with positive probabilities identifies each $v_x(S)$ under consistency, conditional randomization and appropriate interference restrictions. The displayed finite linear map then identifies every $\phi_k(x)$. This is an identification sufficiency statement, not a recommendation that one can literally remove a child's networks or skills.

There is also a randomized-permutation estimator. Draw a channel permutation uniformly, sample a predecessor set $S_k$, and randomize between the adjacent interventions $S_k$ and $S_k\cup\{k\}$. Average their mean earnings difference over permutations. The expectation is precisely $\phi_k$. Its feasibility still depends on defining those channel interventions, and its variance depends on earnings dispersion and assignment probabilities. Observing employers' responses to credential disclosure identifies an information intervention holding the studied skill and network distribution fixed. It need not identify the all-population credential Shapley term because that term averages over eight distinct predecessor sets.

An instrument for completing schooling generally changes the population of compliers and the joint bundle of mechanisms. Even a valid exclusion restriction at the schooling level does not create exclusions for individual channels. Local total effects and modular mechanism effects therefore require separate equations and distinct assumptions.

\section{Transport to a second population}
Let source and target covariate laws be $F_0,F_1$, with $F_1\ll F_0$. Suppose source experiments identify $v_{0,x}(S)$. If all conditional intervention means are invariant, the target mechanism mean is
\[
 \Phi_{1,k}=\int\phi_{0,k}(x)\,dF_1(x)
 =E_0\!\left[\frac{dF_1}{dF_0}(X)\phi_{0,k}(X)\right].
\]
This follows directly by change of measure and linearity. It requires population overlap and mechanism invariance, not merely similar observed earnings distributions.

For an explicit sensitivity calculation assume
$|v_{1,x}(S)-v_{0,x}(S)|\le\delta_S(x)$ for every subset, with integrable bounds. Let $w_S=|S|!(3-|S|)!/4!$. The triangle inequality gives
\[
 |\Phi_{1,k}-\int\phi_{0,k}\,dF_1|
 \le\int\sum_{S\subseteq K\setminus\{k\}}w_S
       [\delta_{S\cup\{k\}}(x)+\delta_S(x)]\,dF_1(x).
\]
For a common bound $\delta$, the right side is $2\delta$ because the weights sum to one. This is a valid coordinate bound, not a characterization of the joint identified set: the same sixteen shifts enter multiple coordinates, and their efficiency relation couples the contributions. Regions with no source overlap require new restrictions or target experiments.

\section{Remaining gap}
The catalogue asks for a credible joint mechanism and heterogeneity analysis in a substantive schooling setting. That requires actual modular intervention definitions, causal evidence about intermediate channel configurations, and defensible population transport assumptions. This note establishes what endpoint schooling evidence cannot do and what stronger experiments would identify. It supplies no earnings estimates, tuition-cost accounting, general-equilibrium credential response, or empirically justified mechanism shares, so the broad target remains unresolved.

\section{Completion Estimate}
50\%

This is a post hoc, heuristic estimate for the full catalogue target.
The attempt proves the interaction content of education Shapley allocations, sharp endpoint-only mechanism nonidentification, and identification and transport sensitivity under modular experiments. Full heterogeneous returns still require feasible channel interventions, observed intermediate configurations, justified target-population overlap, and empirically validated mechanism invariance.

\begin{thebibliography}{9}
\bibitem{ak} J. D. Angrist and A. B. Krueger (1991), \emph{Does Compulsory School Attendance Affect Schooling and Earnings?}, QJE 106(4), 979--1014. Working-paper record and abstract: \url{https://www.nber.org/papers/w3572}. Journal DOI: \url{https://doi.org/10.2307/2937954}.
\bibitem{hlt} J. J. Heckman, L. J. Lochner and P. E. Todd (2006), \emph{Earnings Functions, Rates of Return and Treatment Effects: The Mincer Equation and Beyond}, Handbook of the Economics of Education 1, 307--458. Primary working-paper record: \url{https://www.nber.org/papers/w11544}; published chapter: \url{https://doi.org/10.1016/S1574-0692(06)01007-5}. These references support the distinctions among causal earnings effects, heterogeneity and investment returns; no mechanism-share estimate is imported here.
\end{thebibliography}

\end{document}
